%% Foreign Function Interface
\chapter{Foreign Function Interface}
\label{ch:ffi}

Lambdust provides a memory-safe, capability-based Foreign Function Interface (FFI) for interoperating with C libraries while maintaining type safety and security through static analysis and runtime checks.

\section{FFI Overview}
\label{sec:ffi-overview}

The FFI system provides:

\begin{itemize}
\item \textbf{Type Safety}: Foreign function signatures are statically checked
\item \textbf{Memory Safety}: Automatic marshaling prevents buffer overflows and dangling pointers  
\item \textbf{Capability Security}: Access to foreign functions requires explicit capabilities
\item \textbf{Error Handling}: Foreign function errors are converted to Scheme exceptions
\item \textbf{Garbage Collection Integration}: Foreign pointers are tracked by the GC
\end{itemize}

\section{Foreign Declarations}
\label{sec:foreign-declarations}

\indexentry{foreign declaration}
\indexentry{foreign-declare}

Foreign functions must be declared before use:

\begin{grammar}
\production{\var{foreign-declaration}}{\produces}{(foreign-declare \var{library} \arbno{\var{binding}})}{}
\\
\production{\var{binding}}{\produces}{(\var{name} :: \var{foreign-type})}{}
\\
\production{\var{foreign-type}}{\produces}{(-> \arbno{\var{c-type}} \var{c-type})}{}
\end{grammar}

\begin{example}[Foreign Declarations]
\begin{lambdustcode}
;; Math library functions
(foreign-declare "libm.so"
  (sin :: (-> Float64 Float64))
  (cos :: (-> Float64 Float64))
  (sqrt :: (-> Float64 Float64))
  (pow :: (-> Float64 Float64 Float64)))

;; String functions
(foreign-declare "libc.so" 
  (strlen :: (-> CString Size))
  (strcmp :: (-> CString CString Int32))
  (strdup :: (-> CString CString))
  (free :: (-> Pointer Void)))

;; File I/O
(foreign-declare "libc.so"
  (fopen :: (-> CString CString FilePointer))
  (fclose :: (-> FilePointer Int32))  
  (fread :: (-> Pointer Size Size FilePointer Size))
  (fwrite :: (-> Pointer Size Size FilePointer Size)))

;; System calls
(foreign-declare "libc.so"
  (getpid :: (-> Void Int32))
  (getenv :: (-> CString CString))
  (setenv :: (-> CString CString Int32 Int32)))

;; Custom library
(foreign-declare "./mylib.so"
  (my-function :: (-> Int32 Int32 Int32))
  (init-library :: (-> Void Int32))
  (cleanup-library :: (-> Void Void)))
\end{lambdustcode}
\end{example}

\section{C Type System}
\label{sec:c-types}

\indexentry{C types}
\indexentry{type marshaling}

The FFI supports comprehensive C type mapping:

\section{Primitive Types}
\label{sec:primitive-c-types}

\onecolumn
\begin{table}[H]
\centering
\caption{C Type Mapping}
\label{tab:c-types}
\begin{tabular}{|l|l|l|}
\hline
\textbf{C Type} & \textbf{Lambdust Type} & \textbf{Description} \\
\hline
\code{void} & \code{Void} & No value \\
\code{char} & \code{Int8} & 8-bit signed integer \\
\code{unsigned char} & \code{UInt8} & 8-bit unsigned integer \\
\code{short} & \code{Int16} & 16-bit signed integer \\
\code{unsigned short} & \code{UInt16} & 16-bit unsigned integer \\
\code{int} & \code{Int32} & 32-bit signed integer \\
\code{unsigned int} & \code{UInt32} & 32-bit unsigned integer \\
\code{long long} & \code{Int64} & 64-bit signed integer \\
\code{unsigned long long} & \code{UInt64} & 64-bit unsigned integer \\
\code{float} & \code{Float32} & 32-bit floating-point \\
\code{double} & \code{Float64} & 64-bit floating-point \\
\code{size\_t} & \code{Size} & Architecture-dependent size \\
\code{char*} & \code{CString} & Null-terminated string \\
\code{void*} & \code{Pointer} & Generic pointer \\
\hline
\end{tabular}
\end{table}
\twocolumn

\section{Compound Types}
\label{sec:compound-c-types}

\begin{example}[Compound C Types]
\begin{lambdustcode}
;; Array types
(define-c-type (Array n T) (* T))   ; T[n] in C

;; Structure types
(define-c-struct Point  
  (x :: Float64)
  (y :: Float64))

(define-c-struct Person
  (name :: CString)
  (age :: Int32)
  (height :: Float64))

;; Union types  
(define-c-union Value
  (i :: Int32)
  (f :: Float64)
  (s :: CString))

;; Function pointer types
(define-c-type Callback (-> Int32 Int32))
(define-c-type CompareFunction (-> Pointer Pointer Int32))

;; Opaque pointer types (for library handles)
(define-c-opaque-type FileHandle)
(define-c-opaque-type DatabaseConnection)

;; Declarations using compound types
(foreign-declare "graphics.so"
  (draw-point :: (-> Point Void))
  (distance :: (-> Point Point Float64))
  (make-person :: (-> CString Int32 Float64 Person))
  (print-person :: (-> Person Void)))

(foreign-declare "libc.so"
  (qsort :: (-> Pointer Size Size CompareFunction Void)))
\end{lambdustcode}
\end{example}

\section{Function Calls}
\label{sec:ffi-calls}

\indexentry{foreign-call}
\indexentry{with-capability}

Foreign functions are called through the capability system:

\begin{grammar}
\production{\var{foreign-call}}{\produces}{(foreign-call \var{function-name} \arbno{\var{argument}})}{}
\production{}{\alternative}{(with-capability \var{capability-name} \var{body})}{}
\end{grammar}

\begin{example}[Foreign Function Calls]
\begin{lambdustcode}
;; Basic function calls with capabilities
(with-capability 'math-library
  (lambda ()
    (let ((x 2.0))
      (display "sin(2) = ") 
      (display (foreign-call sin x)) 
      (newline))))

;; String operations  
(with-capability 'libc
  (lambda ()
    (let ((str1 "Hello")
          (str2 "World"))
      (let ((len1 (foreign-call strlen str1))
            (cmp (foreign-call strcmp str1 str2)))
        (display "Length: ") (display len1) (newline)
        (display "Compare: ") (display cmp) (newline)))))

;; File operations with error handling
(define (read-file-contents filename)
  (with-capability 'file-io
    (lambda ()
      (let ((file (foreign-call fopen filename "r")))
        (if (null-pointer? file)
            (error "Cannot open file" filename)
            (begin
              ;; Read file contents
              (let ((buffer (allocate-buffer 1024)))
                (let ((bytes-read (foreign-call fread buffer 1 1024 file)))
                  (foreign-call fclose file)
                  (buffer->string buffer bytes-read)))))))))

;; Callback functions
(define (my-compare-function :: CompareFunction)
  (lambda (a-ptr b-ptr)
    (let ((a (pointer->int32 a-ptr))
          (b (pointer->int32 b-ptr)))
      (cond ((< a b) -1)
            ((> a b) 1)
            (else 0)))))

;; Use callback with foreign function
(with-capability 'libc
  (lambda ()  
    (let ((array (list->c-array '(3 1 4 1 5 9 2 6))))
      (foreign-call qsort array 8 4 my-compare-function)
      (c-array->list array 8))))
\end{lambdustcode}
\end{example}

\section{Memory Management}
\label{sec:ffi-memory}

\indexentry{memory management}
\indexentry{garbage collection}

The FFI provides safe memory management:

\section{Automatic Memory Management}
\label{sec:auto-memory}

\begin{example}[Automatic Memory Management]
\begin{lambdustcode}
;; Automatic string marshaling
(define (safe-string-length str)
  (with-capability 'libc
    (lambda ()
      ;; String automatically converted to C string
      ;; Memory automatically managed
      (foreign-call strlen str))))

;; Automatic cleanup with finalizers
(define (allocate-managed-buffer size)
  (let ((ptr (foreign-call malloc size)))
    (if (null-pointer? ptr)
        (error "Allocation failed")
        (begin
          ;; Register finalizer to free memory
          (register-finalizer! ptr (lambda (p) (foreign-call free p)))
          ptr))))

;; RAII-style resource management
(define-syntax with-file
  (syntax-rules ()
    ((with-file (var filename mode) body ...)
     (let ((var (foreign-call fopen filename mode)))
       (if (null-pointer? var)
           (error "Cannot open file" filename)
           (dynamic-wind
             (lambda () #f)           ; Setup (already done)
             (lambda () body ...)     ; Body
             (lambda ()               ; Cleanup
               (when (not (null-pointer? var))
                 (foreign-call fclose var)))))))))

;; Usage
(with-file (file "data.txt" "r")
  (let ((buffer (allocate-buffer 1024)))
    (foreign-call fread buffer 1 1024 file)
    (buffer->string buffer)))
\end{lambdustcode}
\end{example}

\section{Manual Memory Management}
\label{sec:manual-memory}

\begin{example}[Manual Memory Management]
\begin{lambdustcode}
;; Explicit allocation and deallocation
(define (manual-buffer-operations)
  (with-capability 'libc
    (lambda ()
      (let ((buffer (foreign-call malloc 1024)))
        (if (null-pointer? buffer)
            (error "Allocation failed")
            (begin
              ;; Use buffer
              (foreign-call memset buffer 0 1024)
              (let ((result (process-buffer buffer)))
                ;; Explicit cleanup
                (foreign-call free buffer)
                result)))))))

;; Reference counting for shared resources
(define-struct shared-resource (ptr refcount))

(define (acquire-shared-resource ptr)
  (let ((resource (make-shared-resource ptr 1)))
    (register-finalizer! resource release-shared-resource)
    resource))

(define (retain-shared-resource resource)
  (set-shared-resource-refcount! resource 
    (+ (shared-resource-refcount resource) 1)))

(define (release-shared-resource resource)
  (let ((new-count (- (shared-resource-refcount resource) 1)))
    (set-shared-resource-refcount! resource new-count)
    (when (= new-count 0)
      (foreign-call free (shared-resource-ptr resource)))))

;; Stack-allocated structures
(define-syntax with-stack-struct
  (syntax-rules ()
    ((with-stack-struct (var struct-type) body ...)
     (let ((var (allocate-stack struct-type)))
       body ...))))  ; Automatically freed when scope exits

;; Usage
(with-stack-struct (point Point)
  (set-point-x! point 3.0)
  (set-point-y! point 4.0)
  (foreign-call draw-point point))
\end{lambdustcode}
\end{example}

\section{Capability System}
\label{sec:capabilities}

\indexentry{capability}
\indexentry{security}

The capability system provides fine-grained access control:

\section{Capability Declaration}
\label{sec:capability-declaration}

\begin{example}[Capability Declaration]
\begin{lambdustcode}
;; Define capabilities with permissions
(define-capability math-library
  (libraries "libm.so")
  (functions sin cos tan sqrt pow log exp)
  (memory-limit 0)      ; No heap allocation
  (file-access none))   ; No file system access

(define-capability file-io  
  (libraries "libc.so")
  (functions fopen fclose fread fwrite fseek ftell)
  (memory-limit 1048576)    ; 1MB heap limit
  (file-access read-write)  ; Full file access
  (paths "/tmp/*" "/var/data/*"))  ; Restrict to specific paths

(define-capability network-client
  (libraries "libc.so" "libssl.so")
  (functions socket connect send recv ssl_connect)
  (memory-limit 10485760)   ; 10MB heap limit
  (network-access client)   ; Client connections only
  (domains "*.example.com" "api.service.org"))

(define-capability database
  (libraries "libpq.so")
  (functions PQconnectdb PQexec PQfinish PQgetvalue)
  (memory-limit 52428800)   ; 50MB heap limit
  (network-access client)
  (domains "db.internal.com"))

;; Capability inheritance
(define-capability full-application
  (inherits file-io network-client database)
  (additional-functions my-custom-function))
\end{lambdustcode}
\end{example}

\section{Capability Usage}
\label{sec:capability-usage}

\begin{example}[Capability Usage]
\begin{lambdustcode}
;; Function requiring specific capabilities
(define (process-data-file :: (-> String ~> {FFI[file-io]} String))
  (lambda (filename)
    (with-capability 'file-io
      (lambda ()
        (with-file (file filename "r")
          (let ((buffer (allocate-buffer 4096)))
            (let ((bytes-read (foreign-call fread buffer 1 4096 file)))
              (let ((content (buffer->string buffer bytes-read)))
                (string-upcase content)))))))))

;; Capability-polymorphic function
(define (generic-computation :: (Pi (C : Capability) (-> Unit ~> {FFI[C]} Natural)))
  (lambda (C)
    (lambda ()
      (with-capability C
        (lambda ()
          ;; Can only use functions allowed by capability C
          (computation-using-capability C))))))

;; Capability delegation
(define (delegate-capability child-capability parent-capability operation)
  (with-capability parent-capability
    (lambda ()
      (with-restricted-capability child-capability  ; Subset of parent
        (lambda ()
          (operation))))))

;; Dynamic capability checking
(define (safe-operation capability-name)
  (if (capability-available? capability-name)
      (with-capability capability-name
        (lambda () (perform-operation)))
      (error "Required capability not available" capability-name)))

;; Temporary capability elevation
(define (elevated-operation)
  (with-temporary-capability 'admin-access 300  ; 5 minute timeout
    (lambda ()
      (perform-admin-task))))
\end{lambdustcode}
\end{example}

\section{Error Handling}
\label{sec:ffi-errors}

\indexentry{error handling}
\indexentry{exception mapping}

The FFI converts C errors to Scheme exceptions:

\begin{example}[Error Handling]
\begin{lambdustcode}
;; Automatic error checking
(define (safe-file-operation filename)
  (with-capability 'file-io
    (lambda ()
      (with-ffi-error-handling
        (lambda ()
          (let ((file (foreign-call fopen filename "r")))
            ;; Null pointer automatically converted to exception
            (let ((size (foreign-call file-size file)))
              (foreign-call fclose file)
              size)))))))

;; Custom error mapping
(define-ffi-error-map
  ((null-pointer? fopen) → (file-not-found-error filename))
  ((< 0 malloc) → (out-of-memory-error size))
  ((= -1 socket) → (network-error "Cannot create socket")))

;; Error recovery
(define (robust-network-call)
  (call/cc 
    (lambda (escape)
      (with-exception-handler
        (lambda (e)
          (cond 
            ((network-timeout-error? e)
             (escape 'timeout))
            ((network-connection-error? e)
             (escape 'connection-failed))
            (else (raise e))))
        (lambda ()
          (with-capability 'network-client
            (lambda () 
              (foreign-call network-operation))))))))

;; Resource cleanup on error
(define (safe-resource-operation)
  (let ((resource #f))
    (dynamic-wind
      ;; Acquire resource
      (lambda ()
        (set! resource (foreign-call acquire-resource)))
      ;; Use resource  
      (lambda ()
        (if (null-pointer? resource)
            (error "Resource acquisition failed")
            (foreign-call use-resource resource)))
      ;; Always cleanup
      (lambda ()
        (when (and resource (not (null-pointer? resource)))
          (foreign-call release-resource resource))))))

;; Errno handling
(define (system-call-with-errno operation)
  (with-capability 'system
    (lambda ()
      (foreign-call set-errno 0)  ; Clear errno
      (let ((result (operation)))
        (let ((error-code (foreign-call get-errno)))
          (if (= error-code 0)
              result
              (error "System call failed" 
                     (foreign-call strerror error-code))))))))
\end{lambdustcode}
\end{example}

\section{Threading and Concurrency}
\label{sec:ffi-threading}

\indexentry{FFI threading}
\indexentry{thread safety}

The FFI handles threading considerations:

\begin{example}[FFI Threading]
\begin{lambdustcode}
;; Thread-safe foreign function calls
(define (thread-safe-operation)
  (with-capability 'thread-safe-library
    (lambda ()
      ;; GIL automatically released for blocking calls
      (foreign-call-blocking long-running-c-function))))

;; Thread-local foreign state
(define (thread-local-context)
  (with-thread-local-capability 'database
    (lambda ()
      (let ((conn (foreign-call db-connect)))
        ;; Connection is thread-local
        (foreign-call db-query conn "SELECT * FROM table")))))

;; Async foreign function calls  
(define (async-foreign-call function . args)
  (future-with-capability 'async-library
    (lambda ()
      (apply foreign-call function args))))

;; Usage
(let ((result-future (async-foreign-call expensive-computation 42)))
  ;; Do other work
  (let ((result (await result-future)))
    (display "Result: ") (display result) (newline)))

;; Callback thread safety
(define (thread-safe-callback :: CCallback)
  (lambda (data)
    (with-scheme-thread-context  ; Ensure proper Scheme context
      (lambda ()
        ;; Safe to call Scheme functions from C callback
        (process-callback-data data)))))

;; Register thread-safe callback
(with-capability 'event-library
  (lambda ()
    (foreign-call register-callback thread-safe-callback)))

;; Foreign function call batching
(define (batch-foreign-calls operations)
  (with-capability 'batch-library  
    (lambda ()
      ;; Minimize FFI overhead by batching
      (foreign-call-batch (list->c-array operations)))))
\end{lambdustcode}
\end{example}

\section{Platform Integration}
\label{sec:platform-integration}

\indexentry{platform integration}
\indexentry{native libraries}

The FFI integrates with platform-specific libraries:

\begin{example}[Platform Integration]
\begin{lambdustcode}
;; Platform-specific declarations
(cond-expand
  (linux 
    (foreign-declare "libc.so.6"
      (epoll-create :: (-> Int32 Int32))
      (epoll-ctl :: (-> Int32 Int32 Int32 Pointer Int32))))
  (darwin
    (foreign-declare "libSystem.dylib"
      (kqueue :: (-> Void Int32))
      (kevent :: (-> Int32 Pointer Int32 Pointer Int32 Pointer Int32))))
  (windows
    (foreign-declare "kernel32.dll"
      (CreateEvent :: (-> Pointer Int32 Int32 CString Pointer))
      (WaitForSingleObject :: (-> Pointer Int32 Int32)))))

;; Cross-platform wrapper
(define (create-event-loop)
  (cond-expand
    (linux (with-capability 'system
             (lambda () (foreign-call epoll-create 10))))
    (darwin (with-capability 'system  
              (lambda () (foreign-call kqueue))))
    (windows (with-capability 'system
               (lambda () (foreign-call CreateEvent null-pointer 0 0 null-pointer))))))

;; Dynamic library loading
(define (load-dynamic-library path)
  (with-capability 'dynamic-loading
    (lambda ()
      (let ((handle (foreign-call dlopen path RTLD-LAZY)))
        (if (null-pointer? handle)
            (error "Cannot load library" path (foreign-call dlerror))
            handle)))))

(define (get-dynamic-symbol handle symbol-name)
  (with-capability 'dynamic-loading
    (lambda ()
      (let ((symbol (foreign-call dlsym handle symbol-name)))
        (if (null-pointer? symbol)
            (error "Symbol not found" symbol-name)
            symbol)))))

;; Plugin system
(define (load-plugin plugin-path)
  (let ((handle (load-dynamic-library plugin-path)))
    (let ((init-fn (get-dynamic-symbol handle "plugin_init"))
          (process-fn (get-dynamic-symbol handle "plugin_process"))
          (cleanup-fn (get-dynamic-symbol handle "plugin_cleanup")))
      (make-plugin handle init-fn process-fn cleanup-fn))))

;; System integration
(define (integrate-with-system-service service-name)
  (cond-expand
    (linux (systemd-integration service-name))
    (darwin (launchd-integration service-name))  
    (windows (windows-service-integration service-name))))
\end{lambdustcode}
\end{example}

\section{Performance Considerations}
\label{sec:ffi-performance}

\indexentry{FFI performance}
\indexentry{optimization}

The FFI is optimized for performance:

\begin{example}[FFI Performance]
\begin{lambdustcode}
;; Fast path for simple calls
(define (optimized-math-call x)
  ;; Minimal overhead for simple numeric functions
  (foreign-call-fast sin x))  ; Direct call, no marshaling

;; Bulk operations
(define (process-large-array data)
  (with-capability 'compute-library
    (lambda ()
      ;; Process entire array in C, avoiding per-element overhead
      (let ((c-array (scheme-array->c-array data)))
        (foreign-call process-array-bulk c-array (array-length data))
        (c-array->scheme-array c-array)))))

;; Memory mapping for large data
(define (process-large-file filename)
  (with-capability 'file-io
    (lambda ()
      (let ((mapping (foreign-call mmap-file filename)))
        (foreign-call process-mapped-data mapping)
        (foreign-call munmap mapping)))))

;; Zero-copy string operations  
(define (zero-copy-string-operation str)
  (with-pinned-string str
    (lambda (c-str)
      ;; String memory pinned, passed directly to C
      (foreign-call string-process-in-place c-str))))

;; Inline assembly for critical operations
(define-inline-asm fast-checksum (data :: Pointer) (size :: Size) :: UInt32
  "
  mov %rdi, %rax    ; data pointer
  mov %rsi, %rcx    ; size  
  xor %edx, %edx    ; checksum = 0
  checksum_loop:
    add (%rax), %edx  ; Add word to checksum
    add $4, %rax      ; Next word
    sub $4, %rcx      ; Decrease remaining
    jnz checksum_loop ; Continue if not zero
  mov %edx, %eax    ; Return checksum
  ")

;; SIMD operations via FFI
(define (simd-vector-add v1 v2)
  (with-capability 'simd-library
    (lambda ()
      ;; Use vectorized instructions for array operations
      (foreign-call vector-add-simd v1 v2 (vector-length v1)))))
\end{lambdustcode}
\end{example}

\section{Security Considerations}
\label{sec:ffi-security}

\indexentry{FFI security}
\indexentry{sandboxing}

The FFI includes security measures:

\begin{example}[FFI Security]
\begin{lambdustcode}
;; Sandboxed foreign function execution
(define (sandboxed-computation data)
  (with-sandbox 'computation-sandbox
    (lambda ()
      ;; Limited capabilities, resource limits enforced
      (with-capability 'restricted-compute
        (lambda ()
          (foreign-call untrusted-computation data))))))

;; Input validation
(define (validated-foreign-call function input)
  (validate-ffi-input input function)  ; Static + runtime validation
  (with-capability 'validated-access
    (lambda ()
      (foreign-call function input))))

;; Buffer overflow protection  
(define (safe-buffer-operation buffer size)
  (with-bounds-checked-buffer buffer size
    (lambda (safe-buffer)
      ;; All buffer accesses are bounds-checked
      (foreign-call process-buffer safe-buffer size))))

;; Capability revocation
(define temp-capability 
  (create-temporary-capability 'network-access 300))  ; 5 minutes

;; Revoke capability early if needed
(when (security-breach-detected?)
  (revoke-capability! temp-capability))

;; Audit logging
(with-audit-logging 'foreign-function-calls
  (lambda ()
    (foreign-call sensitive-operation)))
;; All foreign calls logged for security analysis
\end{lambdustcode}
\end{example}

The FFI provides powerful integration with native code while maintaining the safety and security properties that make Lambdust suitable for high-assurance applications.